% --- BLOQUE DE PREÁMBULO UNIVERSAL ---
\documentclass[11pt, letterpaper, twoside]{article}

\usepackage[top=2cm, bottom=2.5cm, left=2cm, right=2cm]{geometry}
\usepackage{fontspec}

\usepackage[spanish, provide=*]{babel}
\babelprovide[import, onchar=ids fonts]{spanish}
\babelprovide[import, onchar=ids fonts]{english}

% Tipografía serif para imitar el Misal Romano
\babelfont{rm}{Times New Roman}

\usepackage{enumitem}
\setlist[itemize]{label=-}

% --- PAQUETES EDITORIALES LITÚRGICOS ---
\usepackage{xcolor}
\usepackage{multicol}
\usepackage{lettrine}
\usepackage{titlesec}
\usepackage{ragged2e}

\definecolor{rubricred}{HTML}{C8102E}

\newcommand{\rubrica}[1]{\textcolor{rubricred}{\small\textit{#1}}}
\newcommand{\vers}{\textcolor{rubricred}{\textbf{V.\ }}}
\newcommand{\resp}{\textcolor{rubricred}{\textbf{R.\ }}}
\newcommand{\cruz}{\textcolor{rubricred}{✠}}

\titleformat{\section}{\Large\bfseries\centering\color{rubricred}}{}{0pt}{\uppercase}
\titleformat{\subsection}{\large\bfseries\centering\color{rubricred}}{}{0pt}{}
\titleformat{\subsubsection}{\normalsize\bfseries\color{rubricred}}{}{0pt}{}

\renewcommand{\LettrineFontHook}{\color{rubricred}}
\setcounter{DefaultLines}{2}
\setlength{\DefaultFindent}{0.2em}
\setlength{\DefaultNindent}{0em}

\setlength{\parindent}{0pt}
\setlength{\parskip}{0.5em}

% --- Configuración de Seguridad para Compilación en Bruto ---
\catcode 95=12
\DeclareRobustCommand{\VAR}[1]{[Dato Dinámico]}
\DeclareRobustCommand{\BLOCK}[1]{}

\begin{document}

\begin{center}
    {\Huge\textbf{Subsidio Litúrgico Diario}\par}
    \vspace{0.3cm}
    {\large\textsc{IV Domingo de Pascua}\par}
    {\small 2026-04-26\par}
\end{center}

\vspace{0.8cm}

\begin{multicols}{2}
\RaggedRight

% ==========================================
% LAUDES
% ==========================================

% ==========================================
% SANTA MISA 
% ==========================================
\section*{SANTA MISA}





\subsection*{GLORIA}
\noindent\rubrica{Se dice el Gloria.} \par\vspace{0.3cm}

\subsection*{ORACIÓN COLECTA}
\lettrine{D}{ios} todopoderoso y eterno, que nos has dado como pastor a tu Hijo Jesucristo...

\subsection*{LITURGIA DE LA PALABRA}
\noindent\rubrica{Hechos 4, 8-12} \par\vspace{0.2cm}
\lettrine{E}{n} aquellos días, Pedro, lleno del Espíritu Santo, dijo: «Jefes del pueblo y ancianos... \par\vspace{0.2cm}
\vers Palabra de Dios. \par
\resp Te alabamos, Señor.

\subsection*{SALMO RESPONSORIAL}
\noindent\rubrica{Salmo 117} \par\vspace{0.1cm}
\resp R. La piedra que desecharon los constructores es ahora la piedra angular. Aleluya. \par\vspace{0.2cm}
\lettrine{R}{.} La piedra que desecharon los constructores es ahora la piedra angular. Aleluya. \\ 
Den gracias al Señor porque es bueno...

\subsection*{SEGUNDA LECTURA}
\noindent\rubrica{1 Juan 3, 1-2} \par\vspace{0.2cm}
\lettrine{Q}{ueridos} hijos: Miren qué amor nos ha tenido el Padre para llamarnos hijos de Dios... \par\vspace{0.2cm}
\vers Palabra de Dios. \par
\resp Te alabamos, Señor.


\subsection*{ACLAMACIÓN ANTES DEL EVANGELIO}
\resp Aleluya, aleluya. \par
Aleluya, aleluya. Yo soy el buen pastor, dice el Señor. \par
\resp Aleluya. \par\vspace{0.3cm}

\subsection*{EVANGELIO}
\noindent\rubrica{Juan 10, 11-18} \par\vspace{0.2cm}
\lettrine{E}{n} aquel tiempo, Jesús dijo a los fariseos: «Yo soy el buen pastor. El buen pastor da la vida por sus ovejas...» \par\vspace{0.2cm}
\vers Palabra del Señor. \par
\resp Gloria a ti, Señor Jesús.


\subsection*{CREDO}
\noindent\rubrica{Se dice el Credo.} \par\vspace{0.3cm}

\subsection*{ORACIÓN DE LOS FIELES}
Sabiendo que el Buen Pastor escucha la voz de su rebaño...

1. Por la Iglesia. \\ 
2. Por la paz. \\ 
3. Por los enfermos. \\ 
4. Por nosotros. \par\vspace{0.3cm}






% ==========================================
% VÍSPERAS
% ==========================================

\end{multicols}
\end{document}